% --- Patron: Malachai, the Cruel Messenger (Curses & Corruption) ---
\subsection{Malachai, the Cruel Messenger (Curses \& Corruption)}

\begin{tcolorbox}[colback=black!5,colframe=red!70!black,title=\textbf{Content Warning}]
Malachai is a patron of curses, corruption, and Faustian bargains. Themes include addiction, loss of self, and the consequences of desperate choices. While the mechanics focus on supernatural afflictions, the fiction may touch on dark subjects. Please use Lines \& Veils to ensure everyone’s comfort at the table.
\end{tcolorbox}

\textit{Lore.} Malachai was once a divine messenger – a herald of painful truths, sent to deliver news that could not be softened. She was not cruel; she was necessary. But necessity wears. In time, she began to soften the truth, offering false hope alongside real curses. For this, the gods bound her in chains of black iron beneath the cursed city of ZahkʼKozahn.

Now she waits. Chained but not contained. Imprisoned but not silent. Her voice reaches across the world, offering power to the weak, hope to the hopeless, and sweet oblivion to those who cannot bear to face the dawn. Vampirism, lycanthropy, and other cursed conditions flow from her influence. Her followers learn that power demands suffering, mastering gifts that devour from within.

Among the peoples of Fate's Edge, Malachai is the patron of the desperate and the damned. In Zakov, she is the \textit{Salt‑Fang}, whose silver coin appears in the mouth of drowned sailors. In the Valewood, she is the \textit{Thorn‑Giver}, invoked by hunters who curse a rival’s bow. In Linn, she is the \textit{Wave‑Mother’s Shadow}, whose gifts turn the tide of battle – at a price. Wherever hope curdles into despair, Malachai whispers.

\begin{quote}
``I offer you everything you desire, and everything you fear. The payment comes later.''
\end{quote}

\noindent\textbf{Domain Focus}
\begin{itemize}
  \item \textbf{Cursed Gifts:} Faustian bargains, power with hidden costs
  \item \textbf{Corruptive Power:} tainted blessings, addictive magic
  \item \textbf{False Salvation:} hope that destroys, beautiful lies
  \item \textbf{Supernatural Afflictions:} vampirism, lycanthropy, and other curses
\end{itemize}

\subsubsection*{Thiasos and Codex (Runekeeper Options)}
A Runekeeper sworn to Malachai does not bind a familiar of purity or light. Their \textbf{Thiasos} is a small creature touched by her corruption: a black rat with one red eye that gnaws at the bars of its cage, a crow with a broken feather that always points toward the nearest corpse, a toad that secretes a toxin causing fever dreams, a small bat that drinks blood from a silver dish, or a moth whose wing dust leaves a rash that itches for a week. This creature requires no food, but it must be given a small sacrifice of pain (a pinprick, a pinch of salt in a wound) each day. If neglected for a week, it becomes \textit{Compromised} (it grows sullen and its aura of despair thickens, imposing –1 die on all Resolve rolls for its bearer until a proper curse is laid).

The \textbf{Codex} of Malachai is never a book of hope. It may be a chain of iron links, each etched with a curse; a box of cracked bells that toll when a Rite is spoken; a scroll of skin from a sinner who died unshriven; or a set of silver nails driven into a board, each nail a promise broken. Because Malachai’s truths are written in pain, the Codex requires \textit{upkeep}: the Runekeeper must spend an hour each week pricking their finger and letting a drop of blood fall onto each nail (or each bell), or the curses stored there become faint (treat as Compromised until re‑blooded).

\subsubsection*{Patron's Gift (Imbuement)}
Once per scene as an action, a Runekeeper of Malachai may touch a held item (weapon, tool, or garment) and whisper a false promise. For the rest of the scene, the item grants \textbf{+1 die} to \textit{Intimidation} or \textit{Deception} rolls, and the first time each scene the bearer would suffer a supernatural curse, they may ignore it (the curse slides off, seeking another target). After the scene, the user marks \textbf{+1 Obligation} as a sliver of Malachai’s hunger lodges in their shadow.

\subsubsection*{Rites of Malachai}

\paragraph{Rite of the Honeyed Binding (Low, 5 XP)} \hfill \\
\emph{Scene; Self or Near/Touch; Standard Push} \\
\textbf{Tags:} \texttt{[CURSE]}, \texttt{[BOOST]}, \texttt{[DEBT]}, \texttt{[FEAR]} \\
\textbf{Materials:} Blood mixed with sweet wine, or a broken chain link. \\
\textbf{Effect:} Choose one:
\begin{itemize}
  \item \textbf{Honeyed Curse:} Target gains \textbf{+2 dice} to their next roll (counts as Triumph). They take 1 Corruption (unresistable).
  \item \textbf{Binding Curse:} Target gains \textbf{+1 die} to physical actions but suffers Fatigue 1 unless they perform a dominance act (beat a weaker creature, drink blood, break something precious).
\end{itemize}
\textbf{Push It:} (Honeyed) Grant \textbf{+3 dice} and auto‑Triumph; inflict 2 Corruption instead. (Binding) Grant \textbf{+2 dice} but Fatigue 2; mark \textbf{1 SB (Spades)} as the curse’s hunger grows. \\
\emph{Requires: Familiar (\textit{Invoke:} 1 Boon).} \\
\textbf{Invoke:} 1 action; mark \textbf{+1 Obligation}.

\paragraph{Rite of the False Dawn (Standard, 8 XP)} \hfill \\
\emph{Scene; Near; Standard Push} \\
\textbf{Tags:} \texttt{[CURSE]}, \texttt{[DEBT]}, \texttt{[FATE]} \\
\textbf{Materials:} Meteoric iron shard blessed by false light. \\
\textbf{Effect:} Bestow a gift with a hidden curse. Target gains a significant advantage (e.g., +1 die to all rolls for one scene) but accrues a 4‑segment \textit{Corruption} clock. Each time they use the gift, tick the clock. When full, they suffer a related curse (GM chooses: festering wounds, beastly urges, or the attention of a Collector). \\
\textbf{Push It:} The gift is more powerful (+2 dice, or extended duration). The clock advances 2 segments immediately; mark \textbf{1 SB (Diamonds)}. \\
\emph{Requires: Familiar + Codex (\textit{Invoke:} 1 Boon).} \\
\textbf{Invoke:} 1 action; mark \textbf{+1 Obligation}.

\paragraph{Rite of the Cruel Transformation (Standard, 9 XP)} \hfill \\
\emph{Scene; Touch; Standard Push} \\
\textbf{Tags:} \texttt{[CURSE]}, \texttt{[TRANSFORM]}, \texttt{[DEBT]}, \texttt{[FEAR]} \\
\textbf{Materials:} Holy water mixed with graveyard dust. \\
\textbf{Effect:} Infect a target with a supernatural condition (vampiric hunger, bestial rage, shadow‑bind). Lasts one scene. Start a 6‑segment \textit{Transformation} clock. Each time the target resists feeding or transformation, tick the clock. When full, the condition becomes permanent. \\
\textbf{Push It:} The condition becomes permanent immediately; target gains \textbf{+1 die} to relevant actions but is marked by Malachai (a Collector will visit within a week). Mark \textbf{1 SB (Hearts)}. \\
\emph{Requires: Familiar + Codex (\textit{Invoke:} 1 Boon).} \\
\textbf{Invoke:} 1 action; mark \textbf{+1 Obligation}.

\paragraph{Rite of the Messenger's Burden (High, 13 XP)} \hfill \\
\emph{Ritual; Self; High Push} \\
\textbf{Tags:} \texttt{[CURSE]}, \texttt{[BOOST]}, \texttt{[PRESSURE]}, \texttt{[DEBT]} \\
\textbf{Materials:} Manacles worn while speaking her name. \\
\textbf{Effect:} Channel Malachai directly. Gain \textbf{+2 dice} to one action per beat this scene, but each use marks 1 SB as her influence spreads. Start an 8‑segment \textit{Corruption's Grip} clock. When it fills, a Collector arrives to demand a memory, a name, or a year of your life. \\
\textbf{Push It:} \textbf{+3 dice} but mark 2 SB, suffer Fatigue 1, and her voice becomes permanent – you hear whispers of cruel bargains in quiet moments. \\
\emph{Requires: Familiar + Codex + Tier III (\textit{Invoke:} 2 Boons).} \\
\textbf{Invoke:} Extended rite; mark \textbf{+2 Obligation}. \\
\emph{Obligation:} 7 segments.

\paragraph{Rite of the Chained Angel's Tribute (High, 14 XP)} \hfill \\
\emph{Extended; Zone; High Push} \\
\textbf{Tags:} \texttt{[CURSE]}, \texttt{[AREA]}, \texttt{[SACRIFICE]}, \texttt{[FATE]} \\
\textbf{Materials:} A chapel bell cracked by divine lightning, or a silver coin that cannot be spent. \\
\textbf{Effect:} Corrupt a sacred space or consecrate a site to Malachai. Within the zone:
\begin{itemize}
  \item Supernatural effects (curses, dark magic) gain \textbf{+1 Effect}.
  \item Blessings and healing spells generate +1 SB on a failure.
  \item Any creature that accepts a gift from Malachai within the zone immediately ticks a \textit{Corruption} clock twice.
\end{itemize}
Start a 10‑segment \textit{Divine Perversion} clock. When it fills, the site becomes a permanent beacon for the Worn – the ghouls of ZahkʼKozahn. \\
\textbf{Push It:} The corruption spreads to adjacent areas; mark \textbf{3 SB (Diamonds)}. \\
\emph{Requires: Familiar + Codex + Tier III (\textit{Invoke:} 2 Boons).} \\
\textbf{Invoke:} Extended rite; mark \textbf{+3 Obligation}. \\
\emph{Obligation:} 8 segments.

\subsubsection*{Cantors \& Cults of Malachai}

\textbf{The Cantor’s Song.} A Cantor of Malachai sings in \textit{harmonies of sweet despair} – melodies that begin like a lullaby and end like a scream. Their voice is used to lure the desperate into bargains, to make a village accept a curse as a blessing, or to sing a monster to sleep so that it may be bound. Malachai’s Cantors are terrifying not because they are evil, but because they are willing to \textit{do evil} – they will poison the well to teach the town to dig deeper, or curse a child to save a hundred adults. They are hated and needed in equal measure.

\textbf{The Cult of the Leashed Hunger.} Malachai’s cults gather in places where suffering is already present: plague hospitals, slums, refugee camps, and the cells of the condemned. The most fanatical cell, the \textit{Leashed Hunger}, meets in catacombs beneath ZahkʼKozahn’s ruins, where they drink water from the Well of Ichor and carve the names of their victims into their own flesh with silver needles. A ritual involves each member naming a curse they have inflicted (or allowed) and a mercy they have buried because of it. Then they prick their finger and let a drop of blood fall into a cracked bell; the bell tolls once, and a Collector appears to witness. Their creed is whispered, not spoken: \textit{“We are the knife that cuts the rot. The wound heals, but the scar remembers us.”} Outsiders are not welcome; they are patients, or potential tools, or future victims. Those who die serving Malachai become the Worn – shambling ghouls that guard the cursed city.

\subsubsection*{Witchcraft of Malachai (Lay / Curse)}

A witch who walks with Malachai is a \textit{curse‑smith} or \textit{affliction‑weaver} – one who crafts pain into purpose. Their magic works through the materials of suffering: a needle that sews a nightmare into a pillow, a forge that tempers a blade to always strike the wielder’s loved ones, or a garden where belladonna and hemlock grow in geometric patterns. Their tools are never hammers; they use needles, threads, and the small iron teeth of cursed combs. In Aeler, a True‑Mason might carve a keystone that makes every tenth step collapse, but a witch of Malachai would instead sew a curse‑seam into the mortar, so that the tunnel itself remembers every injury it has caused.

\textbf{The Witch’s Tool: The Splintered Needle.} A needle of black iron, snapped in half and re‑joined with a drop of the witch’s blood. Once per session, the witch may sew a “curse seam” into a garment, a threshold, or a skin. The curse is small but specific: a button that falls off at the worst moment, a door that sticks when someone is pursued, a hand that trembles when holding a weapon. The target suffers \textbf{–1 die} on a single type of roll (e.g., Stealth, Melee, Sway) for the scene. The witch may unstitch the curse as a free action, but doing so marks 1 Fatigue.

\textbf{A Hedge Gift: Borrowed Scars (2 XP).} Once per scene, the witch may touch a fresh wound (inflicted within the last hour) and speak a single word of Malachai’s name. For the remainder of the scene, the witch may transfer the pain of that wound (as 1 Fatigue) to another creature within Near range. The target resists with Spirit+Resolve (DV 3); on a failure, they suffer the Fatigue and the witch’s original wound heals slightly (clear 1 Fatigue from themselves). On a success, the witch takes the Fatigue anyway – the curse rebounds.

\subsubsection*{Malachai's Corruption Manifestations}

\begin{tblr}{
  colspec = {Q[c,1.8cm] X[2.5] X[2.5]},
  rowhead = 1,
  row{odd} = {bg=gray!10},
  row{1} = {bg=gray!30, font=\bfseries},
  hlines,
}
Level & Benefit & Cost / Quirk \\
1 & \textbf{Sharpened Senses:} +1 die to Notice in dim light. & \textbf{Unsettling Presence:} –1 die to social warmth checks. \\
2 & \textbf{Primal Vigor:} +1 die to Athletics and physical actions. & \textbf{Dark Hunger:} 1 Fatigue if a day passes without feeding (blood or rare herbs). \\
3 & \textbf{Unnatural Resilience:} Once per session, resist one Condition. & \textbf{Beastly Tells:} Failures cause visible physical changes (slitted eyes, fur patches). \\
4 & \textbf{Cursed Vigour:} Natural weapons count as magical; +1 Effect. & \textbf{Transformed Body:} Permanent physical change (eye colour, skin texture, etc.). \\
5 & \textbf{Regenerative Curse:} Heal 1 Fatigue or Condition per rest. & \textbf{Corruption's Price:} Pushing rites adds permanent Corruption (cannot be cleared). \\
6+ & \textbf{Monstrous Ascendant:} +2 dice to Body‑based actions; your unarmed strikes deal +1 Harm and count as magical. & \textbf{Loss of Self:} The GM may compel curse‑driven behaviour (feed, rage, corrupt) once per session. You can resist with Spirit+Resolve (DV 4). On failure, you act according to the curse for one exchange. Your eyes now glow faintly red in darkness – a visible sign of Malachai’s claim. \\
\end{tblr}

\subsection*{Playstyle Notes}
Malachai offers power that corrupts from within. Followers gain immediate benefits but accumulate permanent drawbacks. The corruption progression leads toward becoming a supernatural creature with great power but diminishing humanity. Ideal for players who enjoy tragic characters, Faustian bargains, and exploring the price of power.

\subsection*{Rivalries \& Obligations}

\textbf{Major Rivalries:}
\begin{itemize}
  \item \textbf{The Oath of Flame \& Light:} Direct opposition – sacred vows vs. corrupted promises.
  \item \textbf{The Daughters of the Unbroken Cord:} Both deal in debt, but Daughters harvest generational obligation; Malachai’s followers burn out fast. Collectors may serve both, but a Daughter sees a Malachai addict as a wasted asset. Territory wars over hospices and gambling dens are common.
  \item \textbf{Livaea:} Both use social manipulation, but Malachai destroys while Livaea binds.
\end{itemize}

\textbf{Hard Obligation Triggers:}
\begin{itemize}
  \item \textbf{7+ Segments:} Must corrupt something pure – defile a sacred site, break a true vow, or offer a loved one’s name to a Collector.
  \item \textbf{10+ Segments:} Malachai’s voice becomes overwhelming; refusal risks losing powers. A Collector appears and demands a permanent scar (a visible Mark of Malachai) – refusal means the character falls into the Worn (becomes an NPC ghoul).
\end{itemize}

\noindent\textbf{Emphasises}
\begin{itemize}
  \item \textbf{Cursed Gifts:} power that always costs more than expected
  \item \textbf{Addiction as Magic:} once you start, stopping is harder than any curse
  \item \textbf{False Salvation:} every cure is a waiting room
  \item \textbf{Supernatural Afflictions:} becoming the monster to fight monsters
  \item \textbf{The Chained Angel:} patient, hungry, and never truly contained
\end{itemize}